\part{Mathematical and Computational Modeling}
\label{part:modeling}

This part presents mathematical and computational models of ME/CFS pathophysiology, including:

\begin{itemize}
\item \textbf{Biochemical process models}: Detailed mathematical descriptions of energy metabolism, immune function, and other key processes
\item \textbf{Temporal evolution models}: How symptoms develop and progress over time
\item \textbf{Response to stimuli}: Mathematical modeling of how the body responds to exertion, infections, treatments, and other inputs
\item \textbf{Multi-system integration}: Models connecting different physiological systems
\item \textbf{Predictive models}: Simulating disease trajectories and treatment responses
\item \textbf{Causal hierarchy analysis}: Formal testing of which mechanisms are trigger-capable root causes versus amplifiers versus consequences, using sensitivity analysis, bifurcation theory, and lock removal protocols
\end{itemize}

These models synthesize biological understanding into quantitative frameworks that can generate testable predictions and guide therapeutic strategies. The part concludes with a capstone analysis (Chapter~\ref{ch:causal-hierarchy-formal}) that formally tests the causal hierarchy proposed in Part~II using the integrated 67-variable ODE model.

